Beide Monde umkreisen denselben Zentralkörper, Jupiter. Nach dem dritten Kepler'schen Gesetz ist deshalb
\begin{align}
 \frac{a_\text{G}^3}{T_\text{G}^2} &= \frac{a_\text{I}^3}{T_\text{I}^2}
\end{align}
und die grosse Halbachse der Ganymed-Bahn
\begin{align}
 a_\text{G} &= \aGaF \\
   &= \aIo\cdot\left(\frac{\TGaO}{\TIoO}\right)^{2/3} \\
   &= \aGa \approx \result{\aGaP{3}{2}}
\end{align}
